%! TeX program = lualatex
\documentclass[../main.tex]{subfiles}
\begin{document} \section{The Lotka-Volterra equations}

Systems of differential equations \hlwarn{Edited: ``System of differential equations are'' \(\to\) ``Systems …''} are not always linear and first-order. The example below showcases a non-linear, first-order system --- the Lotka-Volterra \hlwarn{Edited: ``Lokta-Volterra'' \(\to\) ``Lotka-Volterra'' (all occurrences, incl. the section title)} equations.

This following example is adapted from the book \emph{Modelling Life} by Garfinkel, Shevtsov and Guo.

\begin{example} 
  Let's play the role of Fisheries and Oceans Canada (\href{https://www.dfo-mpo.gc.ca}{\ttfamily dfo-mpo.gc.ca}) concerning an ecosystem of sharks (predators) and tuna (prey) \hlwarn{Edited: ``preys'' \(\to\) ``prey''} in the Atlantic provinces.  Fisheries report an increase in shark sightings.  As a result, concerns for a diminishing tuna population and fishing boats safety are mounting.  Calls for relocation of a portion of sharks are picking up momentum. 

  As policy makers, we ask ourselves: \hlmain{Is relocating sharks a good idea?} % Can you see an optimal solution on the vector field?

  \medskip

  We model tuna-shark population dynamics using the following observations.
  \begin{itemize}
    \item The ecosystem is \emph{closed} to migrations.
    \item The per-capita birth rate of tuna is a constant, say \(50\%\).
    \item The per-capita death rate of tuna is a function of the population of sharks, say \(1\%\) of the shark population. \hlwarn{Edited: ``sharks population'' \(\to\) ``shark population'' (both occurrences)}
    \item The per-capita birth rate of sharks is a function of the population of tuna, say \(0.5\%\) of the tuna population.
    \item The per-capita death rate of sharks is a constant, say \(20\%\).
  \end{itemize}

  Here is the vector field for the tuna-sharks system of equations.

  \includestandalone[page=2]{../standalones/plot-vector-field-predatory-prey}

  Lastly, fishing policies should take sustainability and worker safety into account.
  \begin{enumerate}
    \item We cannot fish if the tuna population drops below 20 (thousands).
    \item We cannot fish if the shark population is above 70 (thousands).
  \end{enumerate}

  We can explore the question: \hlmain{Is relocating sharks a good idea?}
\end{example}

\end{document}
